% !PRINTED 378
\weberpage{378}
Demnach liefern die Formeln (13) übereinstimmend

\[
k\,x_4 = \frac{\xi_1}{a_1} + \frac{\xi_2}{a_2} + \frac{\xi_3}{a_3}
\]
oder
\weq{14}{k\,(a_1\,x_1 + a_2\,x_2 + a_3\,x_3) = \frac{\xi_1}{a_1} + \frac{\xi_2}{a_2} + \frac{\xi_3}{a_3}.}
Aus (11) aber folgt noch
\weq{15}{\lambda_1 = -\frac{a_2}{a_3}, \quad \lambda_2 = -\frac{a_3}{a_1}, \quad \lambda_3 = -\frac{a_1}{a_2}.}
Aus (12) ergeben sich dann ferner die drei Relationen
\[
\begin{aligned}
-\frac{\xi_{42}}{\lambda_2} - \frac{\xi_{43}}{\lambda_3} &= \frac{\xi_1}{a_1} + k\,a_1\,x_1 \\
-\frac{\xi_{43}}{\lambda_3} - \frac{\xi_{41}}{\lambda_1} &= \frac{\xi_2}{a_2} + k\,a_2\,x_2 \\
-\frac{\xi_{41}}{\lambda_1} - \frac{\xi_{42}}{\lambda_2} &= \frac{\xi_3}{a_3} + k\,a_3\,x_3.
\end{aligned}
\]
Daraus durch Addition mit Rücksicht auf (14)
\[
\frac{\xi_{41}}{\lambda_1} + \frac{\xi_{42}}{\lambda_2} + \frac{\xi_{43}}{\lambda_3} = -k\,(a_1\,x_1 + a_2\,x_2 + a_3\,x_3),
\]
und folglich
\weq{16}{
\begin{aligned}
\frac{\xi_{41}}{\lambda_1} &= \frac{\xi_1}{a_1} - k\,(a_2\,x_2 + a_3\,x_3) \\
\frac{\xi_{42}}{\lambda_2} &= \frac{\xi_2}{a_2} - k\,(a_3\,x_3 + a_1\,x_1) \\
\frac{\xi_{43}}{\lambda_3} &= \frac{\xi_3}{a_3} - k\,(a_1\,x_1 + a_2\,x_2).
\end{aligned}
}
Unbekannt ist in diesen Formeln noch die Constante $k$. Diese bestimmen wir aus der Bemerkung, dass wir das ganze bisher betrachtete Formelsystem vervierfachen können, indem wir an Stelle von $x_4$ treten lassen $x_4$, $x_5$, $x_6$, $x_7$.

Der Formel (10) entsprechend wollen wir diese vier Functionen linear durch $x_1$, $x_2$, $x_3$ ausdrücken in der Weise:
\weq{17}{
\begin{aligned}
x_4 &= a_{4,1}\,x_1 + a_{4,2}\,x_2 + a_{4,3}\,x_3 \\
x_5 &= a_{5,1}\,x_1 + a_{5,2}\,x_2 + a_{5,3}\,x_3 \\
x_6 &= a_{6,1}\,x_1 + a_{6,2}\,x_2 + a_{6,3}\,x_3 \\
x_7 &= a_{7,1}\,x_1 + a_{7,2}\,x_2 + a_{7,3}\,x_3,
\end{aligned}
}
%CONT>>
